\documentclass{ltjsbook}
\usepackage{docmute} 
\usepackage{../../myStyle} 

\begin{document}
% \chapter{確率的プログラミング；変化点と折線回帰}
\chapter{Probabilistic Programming; Change Points and Piecewise Regression}
\label{x28_ChangePoint}
% さて今回は，今までの線形モデルとはちょっと違うモデリングになります。その名も変化点検出，そして折線回帰です。タイトルだけでも面白そうではありませんか？
Well, this time, we are going to have a slightly different modeling from the linear models we have been working on so far. It's called change point detection, and piecewise regression. Isn't it interesting just by the title?

% さらに今回は次のようなデータを扱います(図\ref{fig:28_01})。何を隠そうこのデータは，私の体重の推移のデータなのです。
Furthermore, we will deal with the following data this time (Figure\ref{fig:28_01}). To tell the truth, this data is about the progression of my weight.
% 私のモーニング・ルーティンとして，朝目が覚めるとトイレに行きます。そして体重計にのり，体重と体脂肪を測定します。その後顔を洗って\footnote{体の内側の水分を出し，表皮に水分をつけないようにすることで，しっかりと肉の量を測ろうとしているのでこの順番になります。}，計測値をiPhoneのアプリに書き込むというのがあります\footnote{ついでにいうと，書き込んだ結果はツイートします。フォロワーの中には私の体重の推移を楽しみにしている人がいるのです。おかしな世の中です。}。このルーチンも早いもので10年近く続けていることになります(一時期やめていたこともありましたが，最古のレコードは2012年06月08日です)。ともかく，このデータを使って，いろいろ遊んでみようと思います。
As part of my morning routine, I go to the bathroom as soon as I wake up. Then, I step on the scale and measure my weight and body fat percentage. After washing my face \footnote {This order is to properly measure the amount of meat by getting the moisture inside the body and not putting moisture on the epidermis.}, I record the measurements in an iPhone app\footnote {By the way, I also tweet the results I wrote down. Among my followers, there are people who look forward to seeing the progress of my weight. It's a funny world. }.  I have been doing this routine for nearly 10 years (there was a time when I stopped, but the oldest record is June 8, 2012). Anyway, I plan to have some fun using this data.


\begin{figure}[ht]
	\begin{center}
		\includegraphics[width=0.8\linewidth,keepaspectratio]{../images/28_ChangePoint/Rplot28_01.png}
		\caption{体重の時系列的推移}
		\label{fig:28_01}
	\end{center}
\end{figure}


% \section{混合分布モデルの応用}
\section{Application of Mixed Distribution Models}
% データが10年分というのはちょっと多いですから，少し時期を区切ることにします。
Because the data spans 10 years, which is a bit much, we will segment it by periods of time.
% 2019年から2020年の2年間に限定したデータが，図\ref{fig:28_02}になります。
The data limited to the two years from 2019 to 2020 is shown in Figure \ref{fig:28_02}.
\begin{figure}[ht]
	\begin{center}
		\includegraphics[width=0.8\linewidth,keepaspectratio]{../images/28_ChangePoint/Rplot28_02.png}
		\caption{2019-2020年のデータ}
		\label{fig:28_02}
	\end{center}
\end{figure}


% これをみると，うーん残念なことに，2020年になると体重が増えてしまっているようですね。2019年は81kg台をうろうろしていたようですが，2020年になると83kg台になっているようです。
Looking at this, unfortunately, it seems that the weight has increased by 2020. It appears that in 2019 the weight was hovering around 81 kg, but it seems to have increased to the 83 kg range in 2020.
% 毎回の体重の計測に偶然の測定誤差がついているとしても，19年と20年とではそもそも体重が違ってしまっているようです。これは平均値の違う正規分布が混合している，混合分布モデルを考えてみることができそうですね。
Even if there are accidental measurement errors in each weight measurement, it seems that the weights in 19th year and 20th year are different in the first place. This suggests that we can consider a mixed distribution model in which normal distributions with different averages are mixed.

% 第\ref{x26_mixed}講を参考に，混合分布モデルを考えてみましょう。筆者の体重には2つの状態があり，軽い方の状態なのか重い方の状態なのかは，確率$\theta$で変わると考えてみます。
Referencing Lecture \ref{x26_mixed}, let's consider a mixed distribution model. My weight has two states, whether it is the lighter state or the heavier state, I assume it changes depending on a probability $\theta$.
% するとデータは，確率$\theta$で$\mu_1$を平均とする正規分布から，確率$1-\theta$で$\mu_2$を平均とする正規分布から得られるわけですから，確率モデルは次のようになります。
Therefore, the data is obtained from a normal distribution with an average of $\mu_1$ at a probability of $\theta$, and from a normal distribution with an average of $\mu_2$ at a probability of $1-\theta$. So, the probability model is as follows.

\[ p(W_j) = \theta_j \times N(\mu_1,\sigma) + (1-\theta_j) \times N(\mu_2,\sigma) \]

ここで$\theta_j$としたのは，観測時点$j$ごとに$\theta$が変わるという想定です。どちらのモードになるのかは一定の数字というより，毎回違っているように思えたのでそのようにしてみました。
されてこをStanで実装するにはどうすれば良いのでしょうか。そう，\verb|log_sum_exp|関数を使って，2つの状態をベクトルで表記し，それを足し合わせる必要があるんでしたね。
実際にコードにしてみたのがコード\ref{stan::28_01}です。

\begin{lstlisting}[language=C++,label={stan::28_01},caption={2つの体重モードモデル}]
data{
  int L;
  array[L] real W;
}

parameters{
  array[L] real<lower=0,upper=1> theta;
  ordered[2] mu;
  real<lower=0> sigma;
}

model{
  for(l in 1:L){
    target += log_sum_exp(
      log(theta[l]) + normal_lpdf(W[l]|mu[1],sigma),
      log1m(theta[l]) + normal_lpdf(W[l]|mu[2],sigma)
    );
  }
  
  mu ~ normal(80,10);
  sigma ~ cauchy(0,5);
}
\end{lstlisting}

これを使って，体重のデータを分析してみましょう。コード\ref{code:28_01}を実行し，図\ref{fig:28_03}のような結果を得ます。
\begin{lstlisting}[language=c++,label={code:28_01},caption={混合分布モデルのコード}]
dat1 <- dat %>%
  dplyr::filter(date > "2019/01/01") %>%
  dplyr::filter(date < "2021/01/01")

model <- cmdstanr::cmdstan_model("changePoint1.stan")
dataSet <- list(L = NROW(dat1), W = dat1$weight)
fit <- model$sample(
  data = dataSet, 
  chains = 4, 
  parallel_chains = 4, 
  seed = 12345)
\end{lstlisting}
\begin{figure}[ht]
   \begin{center}
      \includegraphics[width=0.8\linewidth,keepaspectratio]{../images/28_ChangePoint/Rplot28_03.png}
      \caption{2つの体重のモード}
      \label{fig:28_03}
   \end{center}
\end{figure}

結果の図を見ると，どうやら2019年の初期にも2回ぐらい平均値が高い，「重いモード」が混在しており\footnote{誰がデブモードやねん}，2019年の10月ごろからそちらの比率が増えています。もっとも，2020年の1月にも「軽いモード」に入っている点はあるようですね。

とまあ，このような分析結果になったわけですが，これをみるとデータが82kgぐらいのラインを超えたかどうかで分割されているな，というのがわかります。当然，平均値の違う2つの分布を混ぜたわけですから，平均値の違いによって分かれるわけです。
それにしても2019年の初期はどうしたんでしょうね。1月や4月は重いモード，2-3月と5月以降は軽いモードと，コロコロ入れ替わっています。
ここでは2つのモードのどちらからデータが出てきているのか，ということだけを考えているのでこれでいいのですが，体重というのはそもそも時系列的な変化をするものですから，どこかで重くなった，どこかで軽くなった，というような時系列的なつながりについての情報が，うまくモデル化されていないように思えます。そこで，横軸が時間的な連続であるということに注意して，今度は2021年1月から11月までのデータを見てみましょう。その区間を取り出したのが図\ref{fig:28_04}になります。

\begin{figure}[ht]
   \begin{center}
      \includegraphics[width=0.8\linewidth,keepaspectratio]{../images/28_ChangePoint/Rplot28_04.png}
      \caption{大きく変化したデータ}
      \label{fig:28_04}
   \end{center}
\end{figure}

これは2つのモードがあるというより，途中で大きく変化したというべきではないでしょうか。とくに2021年5月ごろから4kgほど体重がぐぐぐっと下がることがあり，6月以降は下がった体重のレベルに落ち着いているようです。
このように，横軸が時系列だと考えると，5月に何か変化があったのではないか，ということが推察されます。これをモデルで表現してみましょう。

\section{変化点検出}
何かのきっかけでデータの様相がガラリと変わってしまった，その変化したところを変化点と呼び，ここでの課題はその変化点を見つけ出す\keyterm{変化点検出}{へんかてんけんしゅつ}{Change point detection}ということになります。
たとえばセンサーが検出するデータの針が急に違うレベルに変化すると，測定している対象の状態がガラリと変わったのではないか，と考えることができますね。

他にもたとえば，心理学の応用領域として，科学捜査研究所が担当する\keytermL{ポリグラフ検査}{ぽりぐらふけんさ}というのがあります\footnote{皮膚電気活動や心拍，呼吸など複数の整理指標を同時に測定するので\textbf{ポリ}グラフであり，熟練の検査官が反応パターンから被疑者の特別な反応を検出するものです。嘘発見器と呼ばれることがありますが，正確には\keytermL{虚偽検出}{きょぎけんしゅつ}と読んだほうが良いでしょう。最近では\keyterm{隠匿情報検査}{いんとくじょうほうけんさ}{concealed information test}ということもあります。ちなみに科学捜査研究所，通称科捜研は，心理の他にも物理，化学，法医，文書(筆跡鑑定)などの専門領域があり，各都道府県に1つずつ設置されています。}が，これなども針が大きく触れたことで変化をみることになります。

変化点検出はその名の通り「いつ」変化したのかを検出するものです。今回のデータも，だいたい5月ごろに大きな変化があったというのはわかるのですが，いつなのかを特定するのは難しいところです\footnote{とはいえこのデータは筆者自身の体重変化ですので，何があったのかは実はわかっています。5月11日，筆者が帰宅途中に自転車で転倒する事故を起こし，右肩の鎖骨を骨折しました。17日に手術した後，利き手が使いにくいものですから食事の量が減り，結果的に体重(筋肉?)が落ちたというのが真相です。}。これをデータとモデルから明らかにしようというのです。

ベイズ統計はわからないことを確率で表現し，データでその確率情報をアップデートしていくというものです。今回はいつ変化したのがわかりませんから，その変化した時期を$\tau$とし，データ区間の中にその日があると考えて，その区間の一様分布を事前分布とします。その時期$\tau$を境に，データが出てくる分布の位置パラメータがずれると考えるのです。設計図は次のようになります(図\ref{fig:28_05})
\begin{figure}[ht]
   \begin{center}
      \includegraphics[width=0.8\linewidth,keepaspectratio]{../images/28_ChangePoint/slide28_01.png}
      \caption{変化点検出のモデル設計図}
      \label{fig:28_05}
   \end{center}
\end{figure}

設計図をもとに，コードにしてみたのがコード\ref{stan::28_02}になります。
\begin{lstlisting}[language=C++,label={stan::28_02},caption={変化点検出のモデルコード}]
data{
  int L;
  array[L] real W;
}

parameters{
  real<lower=1,upper=L> tau;
  ordered[2] mu;
  real<lower=0> sigma;
}

model{
  for(l in 1:L){
    if(l < tau){
      W[l] ~ normal(mu[2],sigma);
    }else{
      W[l] ~ normal(mu[1],sigma);
    }
  }
  
  tau ~ uniform(1,L);
  mu ~ normal(80,10);
  sigma ~ cauchy(0,5);
}
\end{lstlisting}

日々のデータが順に1からL行目まで並んでいるとします。パラメータ\verb|tau|があり，\verb|l|行目のデータが変化点\verb|tau|より前にあるときは$\mu_2,\sigma$の正規分布から，\verb|tau|より後になれば，$\mu_1,\sigma$の正規分布からデータが出てくると考えるのです。今回は後半の平均値が小さいことが明らかですから，$\mu$のベクトルを\verb|ordered|型で宣言してあります。ベクトル要素の小さい順に並びますから，前半が$\mu_2$，後半が$\mu_1$としています。

これを実行して，変化点がいつになるのかを推定してみましょう。推定結果は出力\ref{MCMCout:out28_01}のようになりました。
\begin{MCMCout}{}{out28_01}
\begin{verbatim}
# A tibble: 4 × 7
  name        EAP       MED       MAP        SD       L95       U95
  <chr> <num:.3!> <num:.3!> <num:.3!> <num:.3!> <num:.3!> <num:.3!>
1 mu[1]    82.770    82.769    82.762     0.062    82.648    82.893
2 mu[2]    78.891    78.892    78.896     0.060    78.775    79.008
3 sigma     0.750     0.749     0.747     0.031     0.693     0.812
4 tau     144.642   144.578   144.431     0.646   143.321   146.190
\end{verbatim}
\end{MCMCout}

データの145行目がちょうどその変化点だといって，ほぼ間違いないようですね。145行目はといいますと，データから5月23日だったことがわかります。
\begin{Rscreen}{いつでしょう}{out28_01}
   \begin{verbatim}
> dat2[145, ]
# A tibble: 1 × 3
  date                weight bodyFat
  <dttm>               <dbl>   <dbl>
1 2021-05-23 06:50:57   80.6      26\end{verbatim}
\end{Rscreen}

違いがわかるようにデータとモデルの推定値を合わせたのが，図\ref{fig:28_06}です。
\begin{figure}[ht]
   \begin{center}
      \includegraphics[width=0.8\linewidth,keepaspectratio]{../images/28_ChangePoint/Rplot28_05.png}
      \caption{検出された変化点と平均値の違い}
      \label{fig:28_06}
   \end{center}
\end{figure}
このようにして，ここで様相が変わったのだなということを，データから見出すことができました。

\section{折線回帰}
先ほどは，変化点を機に平均値が変わる，というモデルでした。では次のような分布の場合はどうなるでしょうか。今度は2016年に注目してみました(図\ref{fig:28_07})。
\begin{figure}[ht]
   \begin{center}
      \includegraphics[width=0.8\linewidth,keepaspectratio]{../images/28_ChangePoint/Rplot28_06.png}
      \caption{2016年の変化}
      \label{fig:28_07}
   \end{center}
\end{figure}
この年はダイエットを志した時期で，最初の半年ぐらいでぐんぐん体重が落ちていってるのがわかります。そして残念なことに，8月ごろでしょうか，ダイエットが終わりリバウンドが始まったのがよくみて取れますね\footnote{このときのダイエット法は，1.ラーメンのスープは飲まない，2.夕食と晩酌は分ける(食べながら飲むのではなく，食べ終わって，お酒だけ飲む),3.お酒を飲むときにおつまみは食べない，という3か条を守るというものでした。お酒が飲みたいので食事の量を減らし，すぐに飲むことにしたので体重が減り，食事の前に飲んでしまえばいいんじゃないかという裏技を見つけてダイエットが崩壊したのを覚えています。}。

さてこのデータに対して，先ほどの変化点検出もできるのですが，変化点の前後で平均値が違う，というような簡単な話ではなさそうです。
変化点の前はどんどん数字が減っていき，変化点のあとはどんどん数字が大きくなっていく，というのが実態です。回帰分析をしたら，変化点前は負の傾きが，変化点後は正の傾きが推定されそうな，そんなデータになっていますね。
ということで，そのイメージができたのであれば，そのままモデルを描いてしまいましょう。設計図が少しごちゃっとしてしまいましたが，ポイントは傾きの係数を前半は負，後半は正に限定しているところでしょうか(図\ref{fig:28_08})。また，変化点は大体このあたり・・・ということでデータの100行目(2016/04/29)から250行目(2016/10/15)の間に置いてみました。
\begin{figure}[ht]
   \begin{center}
      \includegraphics[width=0.8\linewidth,keepaspectratio]{../images/28_ChangePoint/slide28_02.png}
      \caption{折線回帰モデル設計図}
      \label{fig:28_08}
   \end{center}
\end{figure}

設計図をもとに，コードにしてみたのがコード\ref{stan::28_03}になります。
\begin{lstlisting}[language=C++,label={stan::28_03},caption={折線回帰のモデルコード}]
data{
  int L; // data length
  array[L] real W;
  array[L] real X;
}

parameters{
  real<lower=100,upper=250> tau;
  array[2] real beta0;
  real<upper=0> beta1a;
  real<lower=0> beta1b;
  real<lower=0> sigma;
}


model{
  for(l in 1:L){
    if( l < tau ){
      W[l] ~ normal( beta0[1] + (beta1a * X[l]),sigma);
    }else{
      W[l] ~ normal( beta0[2] + (beta1b * X[l]),sigma);
    }
  }
  
  beta0 ~ normal(70,10);
  beta1a ~ normal(0,5);
  beta1b ~ normal(0,5);
  sigma ~ cauchy(0,5);
}
\end{lstlisting}
回帰係数を正・負に限定するところは，パラメータの宣言で\verb|<lower=0>|や\verb|<upper=0>|としていることで表現しています。あとはデータが出てくる正規分布の平均に，線形モデルが入っているだけですね。
これを実行すると，出力\ref{MCMCout:out28_02}のような結果が得られます。データの187行目(2016/08/01)あたりが変化点ですかね。95\%区間で言うと185行目(2016/07/30)から188行目(2016/08/02)の間に変化点があると言えそうです。
\begin{MCMCout}{}{out28_02}
   \begin{verbatim}
# A tibble: 6 × 7
  name           EAP       MED       MAP        SD       L95       U95
  <chr>    <num:.3!> <num:.3!> <num:.3!> <num:.3!> <num:.3!> <num:.3!>
1 beta0[1]    81.871    81.872    81.874     0.075    81.721    82.014
2 beta0[2]    70.548    70.551    70.549     0.380    69.793    71.274
3 beta1a      -0.043    -0.043    -0.043     0.001    -0.044    -0.041
4 beta1b       0.026     0.026     0.026     0.002     0.023     0.029
5 sigma        0.499     0.498     0.497     0.021     0.461     0.543
6 tau        187.135   187.326   187.460     0.832   185.057   188.215
\end{verbatim}
\end{MCMCout}
そして回帰係数をプロットしてみました(図\ref{fig:28_09}。なかなかデータにフィットしていそうです。


\begin{figure}[ht]
   \begin{center}
      \includegraphics[width=0.8\linewidth,keepaspectratio]{../images/28_ChangePoint/Rplot28_07.png}
      \caption{折線回帰モデル推定結果}
      \label{fig:28_09}
   \end{center}
\end{figure}


このままでもよいのですが，変化点と折れたポイントが合致しないのはなんだか気持ち悪いですね。これを合わせることを考えてみたいと思います。
変化点$\tau$のあるところで2つの回帰線は交わる，つまり同じ値になるはずですから，$\beta_{0a}+\beta_{1a}\tau = \beta_{0b}+\beta_{1b}\tau$という式が成り立つはずです。ここから逆算して，
\[ \beta_{0a} + \beta_{1a}\tau - \beta_{1b}\tau = \beta_{0b} \]
% と考えることができますから，推定するパラメータを1つ減らすことができます。
Because we can consider this, we can reduce one parameter that we estimate.
\begin{lstlisting}[language=C++,label={stan::28_04},caption={折線回帰のモデルコード2}]
...(前略)...
parameters{
  real<lower=100,upper=250> tau;
  real beta0a;
  real<upper=0> beta1a;
  real<lower=0> beta1b;
  real<lower=0> sigma;
}

transformed parameters{
  real beta0b;
  beta0b = beta0a + ((beta1a-beta1b) * tau);
}
...(後略)...
\end{lstlisting}

% このコード\ref{stan::28_03}にあるように，\verb|beta0b|を計算式から出すようにして，改めて推定をした結果が図\ref{fig:28_10}になります。これだと変化点が前にずれて，データの169行目(2016/07/10)で心がポッキリ折れていることがわかります。
As shown in this code \ref{stan::28_03}, the result of re-estimation after removing \verb|beta0b| from the calculation formula is shown in the figure \ref{fig:28_10}. With this, we can see that the breakpoint has moved earlier, and we can see that the spirit is definitively broken at the 169th line of data (July 10, 2016).

\begin{figure}[ht]
	\begin{center}
		\includegraphics[width=0.8\linewidth,keepaspectratio]{../images/28_ChangePoint/Rplot28_08.png}
		\caption{変化点を合わせた折線回帰モデル}
		\label{fig:28_10}
	\end{center}
\end{figure}


% このように，モデリングの力を借りるとわからなかった変化点がどこにあるのか，そしてその変化点の前後で予測モデルを変えるといったようなことが簡単に表現できます。ただの回帰分析であっても，データにあったモデルを考えることでさまざまな応用可能性が出てくるのではないでしょうか。
In this way, by using the power of modeling, it becomes easy to express things such as where there is a change point that was not known before, and changing the prediction model before and after that change point. Even it's just a regression analysis, by considering a model that fits the data, various application possibilities may arise.


% ところで，今回は時系列的なデータに対して回帰分析を行いました。この時の説明変数$X_i$は日付，あるいは「何日目か」というデータだったわけです。しかし回帰分析には重要な仮定として，各データ点は(同じ分布から)独立に得られていると言うものがありました。たとえば身長と体重の回帰分析，という話をするときに，各データ点にあたる個々人の身長・体重は，互いに影響し合わない独立だったはずであり，だからこそ尤度の計算の時に各データの尤度を掛け合わせていくことができたのでした。それに対し，体重のデータは明らかに違います。すなわち，今日の体重が明日の体重に影響している，それどころか今日の体重は明日の体重と大いに関わっているはずですし，明日明後日とその後の日々にも少なからず影響しているはずなのです。
By the way, this time we conducted a regression analysis on time-series data. In this case, the explanatory variable $X_i$ was the date, or "how many days". However, a crucial assumption in regression analysis is that each data point is obtained independently (from the same distribution). For example, in discussing the regression analysis of height and weight, each individual's height and weight, which correspond to each data point, were supposed to be independent of each other, and that's why we were able to multiply the likelihood of each data point when calculating the likelihood. In contrast, weight data is clearly different. In other words, today's weight affects tomorrow's weight, let alone, today's weight should be significantly related to tomorrow's weight, and it should be having more or less impact in the following days.

% このように考えると，独立していないデータ点，時系列的なデータに対しては異なる分析をする必要がありそうです。次回はこの時系列的なデータを扱うモデルを紹介していくことになります。
When you think about it this way, it seems necessary to do a different analysis for independent data points and time-series data. Next time, I will introduce a model that handles this time-series data.

% \section{課題}
\section{Assignment}
% 2015年の1月1日から2015年の12月1日の間にも，体重の減少と増加が切り替わったところがありそうです。
It seems there were periods between January 1, 2015, and December 1, 2015, where fluctuations in weight loss and gain occurred.
% 変化点を合わせた折線回帰を実行し，何月何日に変化点があったのか推定するモデルを分析するR/Stanコードを提出してください。結果の解釈などを，スクリプトのコメントアウトや別添ファイルなどで提供してもらえるとなお良いです。もちろんRmdファイルでの提出であれば完璧です。なお提出されたコード単体でバグがなく動くことが確認できないものは，未提出扱いになります。コードの書き方などわからないところがあれば，曜日別TAか小杉までメールで連絡し，指導を受けてください。
Please submit R/Stan code that performs broken line regression with matched change points, and analyzes a model that estimates on what month and day the change points occurred. It would be great if you could provide interpretations of the results, etc., in the comments out of the script or in a separate file. Of course, it would be perfect if you could submit it as an Rmd file. If it is not possible to confirm that the submitted code works without bugs on its own, it will be treated as unsubmitted. If you have any questions about how to write the code, please contact the TA for each day of the week or Kosugi by email and receive guidance.



\end{document}
